\begingroup
\ArithmeticHead{337}{R}{{\scriptsize ZUR THEORIE DER REIHE $F(\alpha,\beta,\gamma,x)$.\quad II.}}
\begin{center}\bfseries BEMERKUNGEN.\end{center}
\par\indent Die in den artt. [1.]---[3] wiedergegebenen Aufzeichnungen folgen in dem Handbuch 18, Bd, unmittelbar auf die Werke III, S. 446---460 abgedruckte Abhandlung, die die Seiten 221---233 des Handbuchs füllt.
Nach der Nr. 140 des \textit{Tagebuchs} ist diese Abhandlung auf Juni 1809 zu datieren; damit befindet sich auch in Übereinstimmung, daß am Ende der S. 220 des Handbuchs (vergl. die Bemerkung von \textsc{Schering}, Werke III, S. 494) die Bemerkung steht: »Geendigt den 28. April 1809.«
\par\indent Unsere artt. [1.]---[3.] sind also um die Mitte des Jahres 1809 niedergeschrieben. --- Die Gleichungen 1), 2), 3) des art. [1.] sind Beziehungen zwischen \textls[120]{verwandten Reihen}*), die den Gleichungen [16], [17], [19] der Abhandlung \textit{circa seriem} Werke III, S. 133 entsprechen.
Die im art. [2.] angegebene Differentialgleichung steht schon auf S. 46, 47 der 1803 begonnenen Scheda Am, wo auch die allgemeine Reihe mit $\alpha,\beta,\gamma$ auftritt; die beiden im art. [2.] vorgenommenen Transformationen finden sich auch Werke III, S. 217 bezw. 224 aus dem Nachlaß abgedruckt.
Der Zweck der im art. [4.] eingeführten Verallgemeinerung erhellt aus dem Beispiel $C=1$, $\gamma=0$.
Zu den numerischen Angaben der artt. [4.] und [5.], die auch aus dem Jahre 1809 stammen dürften, ist folgendes zu bemerken. Den ersten Gleichungen des art. [4.] für $\Pi\,0{,}6$ und $\Pi\,0{,}8$ liegt die Formel
\[\Pi z\,{.}\,\Pi(1-z)=\frac{z(1-z)\pi}{\sin z\pi},\]
\noindent Werke III, S. 148, zu Grunde.
Die mit $\Lambda z$ bezeichnete Größe ist der Logarithmus der im art. 18. der Abhandlung \textit{circa seriem}, Werke III, S. 144 Gleichung [38] erklärten Funktion $\Pi(k,z)$.
Zur der Formel für $\Lambda(z+\omega)$ vergleiche man die Gleichung [60] derselben Abhandlung, a. a. O. S. 153; der Koeffizient von $\omega$ auf der rechten Seite unserer Formel ist nach Gleichung [67], Werke III, S. 154, nichts anderes als $\Psi z$, die Zahl, deren Mantisse in der Handschrift nicht angegeben ist, soll die \textsc{Euler-Mascheroni}sche Konstante sein.
Im art. [5.] bedeutet $\varpi$ die Hälfte der in den Arbeiten über lemniskatische Funktionen ebenso bezeichneten Größe, der dafür angegebene numerische Wert ist der \textsc{Stirling}sche, vergl oben S. 145; von dem hier berechneten Wert von $\log_{10}\varpi$ weicht der Werke III, S. 414 (aus Scheda Aa, S. 3 stammend) angegebene von der 9. Dezimalstelle an ab.
Bei der Berechnung von $\log_{10}\Pi\frac14$ wird die Formel
\[\Pi\frac14=\sqrt[4]{\frac18\pi\varpi^2},\]
\noindent bei der von $\log_{10}\Pi(-\frac14)$ die Formel
\[\Pi(-\frac14)=\Pi\frac14\mathbin{\cdot}\frac{\surd\pi}{\varpi}\]
\noindent benutzt, vergl. Werke III, S. 150. Die Werte von $\log_{10}\Pi\frac14$ und $\log_{10}\Pi\frac34$ sind auf 20 Dezimalstellen auch in der Tafel Werke III, S. 161 enthalten.
Weiterhin kommt allemal die Formel
\[\Pi(z+n)=\Pi z\,{.}\,(z+1)(z+2){.}{.}{.}(z+n)\]
\noindent Werke III, S. 146, Gl. 45 zur Anwendung.
\begin{flushright}\textsc{Schlesinger.}\end{flushright}
\par\Needspace{6\baselineskip}\bigskip\noindent\rule{.25\linewidth}{.4pt}\par\smallskip
{\small\par\indent *) Diese Verdeutschung für \textls[120]{series contiguae} schlägt \textsc{Gauss} selbst in der Anzeige der Abhandlung \textit{circa seriem} vor, siehe Werke III, S. 199.\par}
\ClosingRule
\par\bigskip\noindent{\small\tikz[baseline=0pt,x=1em,y=1em]{\draw[line width=.45pt] (0,.75)--(.5,0); \draw[line width=.45pt] (.25,.375)--(0,0);}\,1.}\hfill 43
\clearpage
\endgroup
